\section{Experience Replay (DQN)}
Los datos consecutivos de un MDP son generalmente dependientes. Para una política fija,
\begin{equation}\label{eq:sequential-dependence}
\begin{aligned}
    P_\pi(s'\mid s)&=\sum_a\pi(a\mid s)P(s'\mid s,a),\\
    \Pr(S_t=s,S_{t+1}=s')&=\Pr(S_t=s)P_\pi(s'\mid s)
    \neq\Pr(S_t=s)\Pr(S_{t+1}=s')\quad\text{en general}.
\end{aligned}
\end{equation}
\begin{itemize}
    \item $P_\pi$: transición inducida por la política; la desigualdad expresa dependencia estadística.
\end{itemize}
La propiedad de Markov elimina la dependencia del historial dado el estado; no hace independientes los estados consecutivos. Además, mientras cambia la política, cambia la distribución de visitas: tampoco hay una distribución idéntica fija~\cite{mnih2013}.

El buffer almacena transiciones y permite muestrearlas aleatoriamente:
\begin{equation}\label{eq:replay}
    \mathcal{D}=\{e_j\}_{j=1}^{N},\qquad
    e_j=(s_j,a_j,r_j,s'_j,d_j),\qquad
    I_1,\ldots,I_m\overset{\mathrm{iid}}{\sim}\operatorname{Unif}\{1,\ldots,N\}.
\end{equation}
\begin{itemize}
    \item $\mathcal{D}$ y $N$: memoria y número de transiciones almacenadas.
    \item $e_j$: experiencia; $d_j=1$ si la transición termina el episodio y cero en otro caso.
    \item $I_1,\ldots,I_m$: índices de un minibatch de tamaño $m$, muestreados con reemplazo de un buffer fijo.
\end{itemize}
Esto reduce la correlación temporal de los minibatches, mezcla experiencias de políticas anteriores y reutiliza datos~\cite{mnih2015}. Condicionado a un buffer fijo, el muestreo con reemplazo es i.i.d. de su distribución empírica; el buffer sigue evolucionando y las experiencias originales siguen relacionadas. Por tanto, replay mejora la optimización, sin garantizar independencia del proceso original ni convergencia. El entrenamiento neuronal también puede usar datos dependientes; i.i.d. es una hipótesis habitual de análisis.
